% Izvor: PDF strana 5, gornji slajd. Redizajn B / etf-v1.
\slajdid{01-s009}
\begin{frame}[label=01-s009]{Logički opis upravljanja liftom}
% element: 01-s009:e01
\Rezultat{Rednim i paralelnim povezivanjem prekidača mogu se razrešiti svi „logički“ problemi pri realizaciji realnih sistema.}
% element: 01-s009:e02
\begin{columns}[T,onlytextwidth]
\begin{column}{.48\textwidth}\Oznaka{Uslovi}\begin{enumerate}
\item Lift je na 3. spratu.
\item Poziv sa 2. sprata „za dole“.
\item Poziv sa 4. sprata „za gore“.
\end{enumerate}\end{column}
\begin{column}{.48\textwidth}\Oznaka{Akcije}\begin{enumerate}
\setcounter{enumi}{3}
\item Pokrenuti lift nagore.
\item Zaustaviti ga na 4. spratu.
\item Otvoriti vrata lifta.
\end{enumerate}\end{column}\end{columns}
% element: 01-s009:e03
\par\vspace{5mm}
\begin{columns}[T,onlytextwidth]
\begin{column}{.43\textwidth}
{\setlength{\parskip}{0pt}Možemo predstaviti promenljivama:\par
\hspace{3mm}\(A\) — položaj na 3. spratu,\par
\hspace{3mm}\(B\) — poziv sa 2. sprata,\par
\hspace{3mm}\(\ldots\)\par
\hspace{3mm}\(Y\) — komanda za kretanje nagore.\par}
\end{column}
\begin{column}{.54\textwidth}
\centering\vspace{12mm}
\(\boxed{Y=f(A,B,\ldots)}\;\longrightarrow\;\text{šema sa prekidačima}\)
\end{column}
\end{columns}

\end{frame}
